% Part: set-theory
% Chapter: ordinals
% Section: introduction

\documentclass[../../../include/open-logic-section]{subfiles}

\begin{document}

\olfileid{sth}{ordinals}{intro}

\olsection{引言}

在 \olref[z][]{chap} 中，我们假设层次结构存在无穷阶段，即 \stagesinf{}（另见无穷公理）。然而，根据 \stagessucc{}，我们不能停留在无穷阶段；我们必须继续下去。因此：在第一个无穷阶段之后的下一阶段，我们形成由第一个无穷阶段中一切可用集合构成的所有可能的汇集；然后重复；再重复；再重复；……

这里暗含的是，我们开始诉诸一种“直观的”数的概念，根据这一概念，在所有自然数\emph{之后}依然可以存在数。具体而言，所涉及的概念正是\emph{超限序数}。本章的目的即是将这一想法严格化。我们将探讨序数的一般概念，然后显式地定义某些集合作为我们的序数。

\end{document}
